% Part: second-order-logic
% Chapter: syntax-and-semantics
% Section: semantic-notions

\documentclass[../../../include/open-logic-section]{subfiles}

\begin{document}

\olfileid{sol}{syn}{sem}
\olsection{المفاهيم الدلالية}

\begin{explain}
تبقى تعريفات \emph{الصحة} و\emph{الاستلزام} و\emph{قابلية الإرضاء}،
وهي أصول البحث المنطقي في الدلالة، على حالها عند الانتقال من الرتبة
الأولى إلى الثانية. وإنما نبنيها الآن على علاقة إرضاء !!{formula}s
الرتبة الثانية. وأما !!{sentence} هذه الرتبة فهي
!!a{formula} قُيّدت جميع متغيراتها، بما فيها متغيرات المحمولات
والدوال.
\end{explain}

\begin{defn}[الصحة]
الجملة~$!A$ \emph{صحيحة منطقيًا}، ورمز ذلك~$\Entails !A$، إذا
وفقط إذا كان~$\Sat{M}{!A}$ في كل~!!{structure}~$\Struct M$.
\end{defn}

\begin{defn}[الاستلزام]
نقول إن مجموعة الجمل~$\Gamma$ \emph{تستلزم} الجملة~$!A$، ونكتب
$\Gamma \Entails !A$، إذا وفقط إذا كانت كل~!!{structure}~$\Struct M$
يتحقق فيها~$\Sat{M}{\Gamma}$ يتحقق فيها كذلك~$\Sat{M}{!A}$.
\end{defn}

\begin{defn}[قابلية الإرضاء]
مجموعة الجمل~$\Gamma$ \emph{قابلة للإرضاء} متى كان
$\Sat{M}{\Gamma}$ في بعض~!!{structure}~$\Struct M$.
فإن لم تكن~$\Gamma$ كذلك قلنا إنها \emph{غير قابلة للإرضاء}.
\end{defn}

\end{document}
